\section*{अध्याय \ref{ch:g12:randvar} --- यादृच्छिक चर और द्विपद बंटन}

\begin{solution}{exo:g12:randvar:1}
$\E(X) = -0.3 + 0 + 0.8 + 0.5 = 1$। $\E(X^2) = 0.3 \times 1 + 0 + 0.4\times4 + 0.1\times25 = 4.4$, इसलिए कोनिग--होयगेंस से $\V(X) = 4.4 - 1 = 3.4$ और $\sigma(X) = \sqrt{3.4} \approx 1.84$।
\end{solution}

\begin{solution}{exo:g12:randvar:2}
\emph{1.} $10$ प्रश्न $p = \frac14$ वाले \omterm{def:g12:condprob:indep}{स्वतंत्र} \omterm{def:g12:randvar:bernoulli}{बर्नूली परीक्षण} हैं: $X \sim \mathcal B\left(10, \frac14\right)$, $\E(X) = 2.5$, $\sigma(X) = \sqrt{10 \times \frac14 \times \frac34} = \sqrt{1.875}
\approx 1.37$।

\emph{2.} $\P(X = 0) = \left(\frac34\right)^{10} \approx 0.056$;
\[
\P(X = 5) = \binom{10}{5}\left(\frac14\right)^5\left(\frac34\right)^5
\approx 0.058;
\quad
\P(X \geq 1) = 1 - \left(\tfrac34\right)^{10} \approx 0.944 .
\]
\end{solution}

\begin{solution}{exo:g12:randvar:3}
\omterm{def:g12:condprob:indep}{स्वतंत्र} और एक जैसे परीक्षण: $X \sim \mathcal B(400,\ 0.05)$, इसलिए $\E(X) = 20$ दावे और
\[
\sigma(X) = \sqrt{400 \times 0.05 \times 0.95} = \sqrt{19} \approx 4.4 .
\]
\end{solution}

\begin{solution}{exo:g12:randvar:4}
प्राप्त भुगतान $P$ $\E(P) = \frac{5 + 6}{6} = \frac{11}{6}$ को संतुष्ट करता है, इसलिए न्यायसंगत दाम $c = \frac{11}{6} \approx 1.83$ यूरो
है। न्यायसंगत लाभ $G = P - \frac{11}6$ \omterm{def:g10:proba:distribution}{प्रायिकताओं} $\frac46, \frac16, \frac16$ के साथ मान $-\frac{11}{6}, \frac{19}{6}, \frac{25}{6}$ लेता है:
\[
\V(G) = \E(G^2)
= \frac{4 \times 121 + 361 + 625}{6 \times 36} = \frac{1470}{216}
= \frac{245}{36} \approx 6.8,
\qquad \sigma(G) \approx 2.6 .
\]
न्यायसंगत खेल का \emph{औसत} लाभ शून्य होता है, पर उसके परिणाम फिर भी
उतार-चढ़ाव करते हैं: न्यायसंगत का अर्थ जोखिम-रहित नहीं है।
\end{solution}

\begin{solution}{exo:g12:randvar:5}
$X \sim \mathcal B(8,\ 0.7)$।
\[
\P(X = 6) = \binom86 (0.7)^6(0.3)^2 \approx 0.296 ;
\]
\[
\P(X \geq 6) = \P(6) + \P(7) + \P(8)
\approx 0.296 + 8(0.7)^7(0.3) + (0.7)^8
\approx 0.296 + 0.198 + 0.058 = 0.552 ;
\]
$\P(X \geq 1) = 1 - (0.3)^8 \approx 0.99993$. बहुलक: $(n+1)p = 6.3$, इसलिए सबसे संभावित मान $k^* = 6$ है (\cref{exo:g12:randvar:7})।
\end{solution}

\begin{solution}{exo:g12:randvar:6}
पहुँचने वाले यात्रियों की संख्या $X \sim \mathcal B(104,\ 0.9)$ है; उड़ान तब अधिविक्रीत हो जाती है जब
$X \geq 101$:
\[
\P(X \geq 101) = \sum_{k=101}^{104}\binom{104}{k}(0.9)^k(0.1)^{104-k}
\approx 0.006 .
\]
सीटों से $4\%$ अधिक टिकट बेचने पर केवल लगभग $0.6\%$ उड़ानों में गड़बड़ होती है ---
यही अधिविक्रय के पीछे का अर्थशास्त्र है।
\end{solution}

\begin{solution}{exo:g12:randvar:7}
\[
\frac{\P(X = k+1)}{\P(X = k)}
= \frac{\binom{n}{k+1}}{\binom nk}\cdot\frac{p}{1-p}
= \frac{n - k}{k + 1}\cdot\frac{p}{1-p},
\]
जहाँ $\binom{n}{k+1} = \binom nk \frac{n-k}{k+1}$ का उपयोग हुआ है। यह अनुपात $\geq 1$ ठीक तब है जब $(n-k)p \geq (k+1)(1-p)$, यानी जब $np - k p \geq k - kp + 1 - p$,
यानी जब $k \leq (n+1)p - 1$। इसलिए \omterm{def:g10:proba:distribution}{प्रायिकताएँ} तब तक कड़ाई से बढ़ती हैं जब तक $k + 1 \leq (n+1)p$ हो और उसके
बाद घटती हैं: \omterm{def:g10:functions:extrema}{अधिकतम} $k^* = \floor{(n+1)p}$ पर प्राप्त होता है (और जब $(n+1)p$ \omterm{def:g10:numbers:sets}{पूर्णांक} हो तब $k^* - 1$
के साथ साझा)।
\end{solution}

\begin{solution}{exo:g12:randvar:8}
\emph{1.} पहली चित उछाल $k$ पर आने का अर्थ है $k-1$ पट फिर चित: \omterm{def:g10:proba:distribution}{प्रायिकता}
$\left(\frac12\right)^{k-1}\cdot\frac12 = 2^{-k}$, $1 \leq k \leq 10$ के लिए। जीतने की कुल \omterm{def:g10:proba:distribution}{प्रायिकता} $\sum_{k=1}^{10} 2^{-k} = 1 - 2^{-10} = \frac{1023}{1024}$ (खेल केवल लगातार दस पटों पर
हारा जाता है)।

\emph{2.} प्रत्याशित भुगतान:
\[
\sum_{k=1}^{10} 2^k \cdot 2^{-k} = \sum_{k=1}^{10} 1 = 10 \text{ यूरो}.
\]
सीमा हटा दें तो योग $\sum_{k\geq1} 1$ अपसरित हो जाता है: प्रत्याशित भुगतान अनंत हो जाता है,
यद्यपि खेल लगभग हमेशा छोटी-सी राशि ही देता है --- यही प्रसिद्ध \emph{सेंट
पीटर्सबर्ग विरोधाभास} है, जो दिखाता है कि अकेली \omterm{def:g12:randvar:exp}{प्रत्याशा} किसी खेल का मूल्य
नहीं नाप सकती।
\end{solution}

\begin{solution}{pb:g12:randvar:1}
\textbf{1.} $X \sim \mathcal B\!\left(3, \frac16\right)$: $\P(X = 0) = \frac{125}{216}$, $\P(X = 1) = \frac{75}{216}$, $\P(X = 2) = \frac{15}{216}$, $\P(X = 3) = \frac{1}{216}$।

\textbf{2.} $\E(X) = 3 \times \frac16 = \frac12$; $V(X) = 3 \times \frac16 \times \frac56 = \frac{5}{12}$।

\textbf{3.} $1$-यूरो के दाँव के सामने प्रत्याशित भुगतान $2\,\E(X) = 1$ यूरो: लाभ $0$
--- ठीक न्यायसंगत (जो विरला है)।

\textbf{4.} बिना वापसी के निकाले परावलंबी होते हैं (दूसरे पत्ते की संभावनाएँ
पहले पर निर्भर करती हैं): द्विपद नहीं --- जाँच-सूची का \omterm{def:g12:condprob:indep}{स्वतंत्रता} वाला ख़ाना
ख़ाली रह जाता है। वापसी सहित: नियत $n = 5$, अचर $p = \frac14$, \omterm{def:g12:condprob:indep}{स्वतंत्र} निकाले: द्विपद।

\textbf{5.} $\E(X) = 6$, $\sigma(X) = \sqrt{4.2} \approx
2.05$। $\E(S) = 5 \times 6 - 10 = 20$; $\sigma(S) = 5\sigma(X) \approx 10.25$।

\textbf{6.} नियत $n = 105$ परीक्षण (टिकट), हर एक में एक ही $p = 0.9$ के साथ
पहुँचना/न-पहुँचना, और \omterm{def:g12:condprob:indep}{स्वतंत्रता} मान ली गई: द्विपद। और स्वीकारोक्ति: परिवार
\emph{एक साथ} उड़ान चूकते हैं और तूफ़ान पूरे विमान ख़ाली कर देते हैं ---
\omterm{def:g12:condprob:indep}{स्वतंत्रता} प्रतिरूप की आस्था की छलाँग है, और सहसंबंधित न-पहुँचना पूँछों को
द्विपद के वादे से मोटा कर देता है।

\textbf{7.} $\E(X) = 94.5$; $\sigma(X) = \sqrt{105 \times 0.9 \times 0.1} =
\sqrt{9.45} \approx 3.07$।

\textbf{8.} $\P(X \geq 101) = \sum_{k=101}^{105}
\binom{105}{k} 0.9^k\, 0.1^{105 - k}$।

\textbf{9.} $\approx 0.0167$: लगभग साठ में एक उड़ान पर किसी न किसी को वंचित होना पड़ता है।

\textbf{10.} प्रति उड़ान $\E(\text{वंचित}) \approx 0.023$ यात्री: प्रत्याशित मुआवज़ा $\approx 19$ यूरो --- और
सामने $1\,000$ यूरो की अतिरिक्त आय। पाँच का अधिविक्रय भारी लाभदायक है; इसीलिए हर
विमान-कंपनी यह करती है।

\textbf{11.} पूँछ-प्रायिकताओं की सारणी से: $n = 105$: $1.7\,\%$; $n = 106$: $4.0\,\%$; $n = 107$:
$8.1\,\%$। नियमों पर खरी उतरने वाली सबसे बड़ी बिक्री $n = 106$ टिकट है।

\textbf{12.} $z = \frac{100.5 - 94.5}{3.07} \approx 1.95$: सीमा \omterm{def:g11:stat:mean}{माध्य} से दो \omterm{def:g11:stat:variance}{मानक विचलन} ऊपर बैठी है, और घंटी-वक्र का
मोटा नियम (``ऊपरी $2\sigma$ पूँछ $\approx 2.5\,\%$'') यथार्थ $1.7\,\%$ के पास आ गिरता है --- द्विपद
के पीछे छाया की तरह चलता चिकना वक्र ही आने वाले अध्यायों का नायक है।

\textbf{13.} $0.98^{20} + 20 \times 0.02 \times 0.98^{19}
\approx 0.94$: ईमानदार खेप $94\,\%$ बार पास हो जाती है।

\textbf{14.} $0.9^{20} + 20 \times 0.1 \times 0.9^{19}
\approx 0.39$: ख़राब खेप फिर भी $39\,\%$ बार निकल जाती है।

\textbf{15.} उत्पादक का जोखिम: अच्छी खेप का अस्वीकार ($\approx 6\,\%$); उपभोक्ता का
जोखिम: ख़राब खेप का स्वीकार ($\approx 39\,\%$)। बड़ा \omterm{def:g10:proba:sample}{प्रतिदर्श} (और उसके अनुपात में
स्वीकृति-सीमा) दोनों को सिकोड़ देता है --- और क़ीमत अधिक निरीक्षण की होती है:
गुणवत्ता की भी एक बजट-रेखा है।

\textbf{16.} $V\!\left(\frac Xn\right) = \frac{V(X)}{n^2} =
\frac{np(1-p)}{n^2} = \frac{p(1-p)}{n}$: \omterm{def:g11:stat:variance}{मानक विचलन} $\sqrt{\frac{p(1-p)}{n}}$। \omterm{def:g10:proba:sample}{प्रतिदर्श} चौगुना कीजिए, शोर आधा हो
जाएगा: कक्षा 10 के उतार-चढ़ाव \omterm{def:g10:numbers:interval}{अंतरालों} का वही $\frac{1}{\sqrt n}$ वाला नियम, अब निकाला हुआ।

\textbf{17.} सीमा भुगतान को $2^{10} = 1024$ यूरो पर बाँध देती है, इसलिए दसों दौरों में
से हर एक \omterm{def:g12:randvar:exp}{प्रत्याशा} में ठीक $1$ यूरो जोड़ता है: $\E = 10$। बिना सीमा वाले खेल की
अनंत \omterm{def:g12:randvar:exp}{प्रत्याशा} पूरी तरह खगोलीय रूप से विरले और खगोलीय रूप से बड़े भुगतानों से
आती थी; और सीमा लगाना यह स्वीकार कर लेना है कि कोई बैंक $2^{50}$ यूरो नहीं देता।
सीमा-बद्ध खेल के लिए न्यायसंगत टिकट-दाम: $10$ यूरो --- और अचानक किसी को कोई
विरोधाभास नहीं रहता।

\textbf{18.} प्रति टिकट प्रत्याशित पुरस्कार-राशि: $2$-यूरो के टिकट के सामने
$\frac{100 + 2 \times 50}{200} = 1$ यूरो: लाभांश $50\,\%$ --- यानी रूले के $2.7\,\%$ का अठारह गुना। लॉटरी इसलिए बची
रहती है कि उसके खिलाड़ी जान-बूझकर दान कर रहे होते हैं: वह ``हानि'' ही तो मक़सद
है।

\textbf{19.} प्रत्याशित दावे: $1\,000 \times 0.01 \times
10\,000 = 100\,000$; प्रीमियम $130\,000$: प्रत्याशित लाभ $30\,000$
यूरो। कुल दावों का \omterm{def:g11:stat:variance}{मानक विचलन}: $10\,000 \times \sqrt{1\,000 \times 0.01 \times 0.99}
\approx 31\,500$ यूरो --- एक साधारण बुरा वर्ष ही पूरा
प्रत्याशित लाभ चट कर जाता है। इसीलिए पूँजी-कोष, पुनर्बीमा, और एक हज़ार
बीमा-पत्रों से कहीं बड़े संचय: बीमाकर्ता बृहत् संख्याओं के नियम पर जीते हैं और
उसकी धीमी चाल के विरुद्ध कोष रखते हैं।

\textbf{20.} पहले जाँच-सूची --- नियत $n$, वही $p$, और \omterm{def:g12:condprob:indep}{स्वतंत्रता} का
\emph{दावा} --- वरना द्विपद है ही नहीं। फिर $\E \pm \sigma$ से दिशा तय कीजिए: \omterm{def:g11:stat:mean}{माध्य} जगह
बताते हैं, विचलन चेतावनी देते हैं। फिर पूँछों का दाम लगाइए: वंचन, ख़राब खेप और
विनाशकारी वर्ष, सब $2\sigma$ से आगे रहते हैं। चेतावनियाँ: \omterm{def:g12:condprob:indep}{स्वतंत्रता} प्रतिरूपकार का
वादा है, संसार का नहीं; और \omterm{def:g12:randvar:exp}{प्रत्याशा} ठीक उसी को अनदेखा कर देती है जो आपको
मिटा देता है। नियम-पुस्तिका के छूटे हुए पन्ने --- आपेक्षिक बारंबारताएँ कितनी
तेज़ी से जमती हैं, और उतार-चढ़ाव की आकृति क्या है --- अगले दो अध्याय हैं।
\end{solution}
